\section{Controlled experiments at a prescribed memory profile}
\label{sec:controlled-memory}
An unrestricted posterior is a useful information state for a proof, but
it is not a free runtime state. We first give a resource-exact dynamic
formulation and then solve a controlled family in which replacing that
formulation by a policy-tree mixture changes the value strictly.

\subsection{A parameterwise occupation state}
Fix a finite horizon $T$, a compact metric parameter space $\Theta$,
finite hidden execution spaces $H_t$, and register alphabets $S_t$ of
sizes at most $K_t$. The visible inputs to each stochastic row and all
legal actions are specified. Any input that must survive a cut belongs
to $S_t$: this includes a chosen action, an unread report buffer, a stop
flag, a selected validation event and a validation accumulator.
One may subdivide an operation into several cuts to expose these costs.
The scheduled time is available only in the clocked convention.

For now rows at different scheduled times are independently prescribed.
Let $\mathcal Q_t$ be the finite product of legal row simplices at stage
$t$. A prescription $q\in\mathcal Q_t$ is a table, chosen at design
time, whose runtime arguments are exactly the declared visible inputs.
After summing the fixed environmental kernels its one-stage transition
is a stochastic matrix $M^q_{t,\theta}$ on the hidden-state/register
pairs. We assume this matrix is jointly continuous in $(q,\theta)$.
Both action selection and state updating occur in this matrix, in their
specified causal order; it is not an arbitrary matrix chosen separately
for each hidden state or each parameter.

For a continuous bundle of probability vectors
$x=(x_\theta(h,s))_{\theta\in\Theta}$ define
\[
 (\mathcal T_{t,q}x)_\theta(h',s')
  =\sum_{h,s}x_\theta(h,s)
              M^q_{t,\theta}((h,s),(h',s')).
\]
The initial bundle $x_0$ is fixed by preparation. Terminal rewards
$r_\theta(h,s)\in[-1,1]$ are continuous in $\theta$.
Absorbing stopped states and frozen validation phases are included in
this execution, not added afterwards without charge. Put
\begin{align}
 B_T(x)&=\min_{\theta\in\Theta}\sum_{h,s}x_\theta(h,s)r_\theta(h,s),
                                                    \label{eq:memory-terminal}\\
 B_t(x)&=\max_{q\in\mathcal Q_t}B_{t+1}(\mathcal T_{t,q}x),
       \qquad 0\le t<T.                              \label{eq:memory-dp}
\end{align}
The argument $x$ is an \emph{offline occupation bundle}, not the
realized history, not an observed parameter, and not an additional
register made available to the machine.

\begin{theorem}[Resource-exact dynamic value and supersolutions]
\label{thm:memory-dp}
For the specified clocked profile and observable-row interface,
$B_0(x_0)$ is the exact maximum worst-case reward over all stochastic
controllers. The maximum is attained by a sequence of legal common
prescriptions. With the norm
$\|x-x'\|=\sup_\theta\sum_{h,s}|x_\theta(h,s)-x'_\theta(h,s)|$,
each $B_t$ is $1$-Lipschitz. Moreover,
\begin{equation}\label{eq:memory-super}
 B_0(x_0)=\inf_H H_0(x_0),
\end{equation}
where the infimum ranges over families of functions on occupation
bundles satisfying
\[
 H_T\ge B_T,\qquad
 H_t(x)\ge H_{t+1}(\mathcal T_{t,q}x)
                       \quad(q\in\mathcal Q_t).
\]
The infimum is attained by $H=B$. No convexification over retained
controllers and no new adversarial choice after a report occurs in
these identities.
\end{theorem}
\begin{proof}
For a fixed sequence $q_0,\ldots,q_{T-1}$, iteration of the displayed
transition gives its actual terminal law for every fixed $\theta$.
The terminal minimum is therefore exactly its worst-case reward.
A runtime randomized controller in this interface is precisely such a
sequence of common conditional rows; independence of fresh randomization
and the declared register specify all its conditional choices.
Persistent random correlations are allowed only through that register.

The occupation bundle at the next stage is determined by the current
bundle and the prescription, without drawing a realized history in the
design problem. Hence optimizing over the remaining prescription
sequence gives \eqref{eq:memory-dp} by ordinary backward induction.
This does not select a different continuation for each unrecorded
history. Each maximizing $q$ is still a single legal table common to
all hidden histories and parameters. At $x_0$ one follows the maximizing
sequence of design prescriptions, computing their occupation bundles
offline; the runtime controller only reads its allowed inputs.

Stochastic matrices contract the $\ell^1$ norm. The terminal minimum is
$1$-Lipschitz in the displayed norm. Induction, using the same $q$ when
comparing two bundles and then taking the two maxima, proves the
Lipschitz assertion. Joint continuity and compactness of $\mathcal Q_t$
give attainment at each stage, and compactness of $\Theta$ gives the
terminal minimum. Any family $H$ in \eqref{eq:memory-super} bounds the
terminal reward along every prescription sequence. Backward induction
therefore gives $H_t\ge B_t$. Equality is admissible, proving
\eqref{eq:memory-super} and its attainment.
\end{proof}

This is a nonlinear supersolution dual, not a least-favorable-prior
saddle theorem at fixed width. Its optimization can be large and
nonconvex. The principle of dynamic programming is classical; the point
of this formulation is that the common-row and memory constraints
remain in the maximizations. If rows are shared across scheduled times,
as for an autonomous controller, choose the shared rows once and hold
them fixed in all subsequent transitions; equivalently augment the
\emph{design} problem by that immutable prescription. Maximizing a
shared row anew at each time would solve a different problem.
A finite union of allowed interfaces can be handled by taking their
maximum, or by a fully wired interface with the prescribed profile and
its legal observable inputs. Neither operation enlarges a runtime
register invisibly.

\subsection{An exact informative-control bottleneck}
Fix integers $m\ge2$ and $1\le K\le m$. The fixed parameter is
$\theta=(i,j)\in[m]\times\{0,1\}$. The first training experiment $R$
reports $i$. Across the next cut the controller retains at most $K$
labels, with an arbitrary stochastic encoder $E(s\mid i)$. It then
chooses one of the $m$ sensing actions with probabilities $D(a\mid s)$.
The selected action is kept in a charged $m$-state action buffer and the
old label is discarded. Action $a$ reports $j$ when $a=i$ and reports
$\perp$ otherwise. The first report is no longer available.

The candidate validation law on $\{0,1\}$ is $\delta_j$; the target is
uniform. An event is frozen before one fresh candidate and one fresh
target validation. Stopping before either training action is permitted.
No independent persistent selector is given to the controller. The
observation symbols are atomic; a bitwise input implementation would
have additional cuts. Denote the maximum worst-case score by $V_{m,K}$.

\begin{theorem}[Exact finite-memory control and its convexification gap]
\label{thm:routing-gap}
In this protocol,
\begin{equation}\label{eq:routing-value}
 V_{m,K}=\frac{1}{2\lceil m/K\rceil}.
\end{equation}
The value obtained by allowing an uncharged ex ante mixture of complete
$K$-label controllers is instead
\begin{equation}\label{eq:routing-mixed}
                         \overline V_{m,K}=\frac{K}{2m}.
\end{equation}
The inequality is strict whenever $K$ does not divide $m$. The matching
controller for \eqref{eq:routing-value} has charged cut profile bounded
by $(1,K,m,3,6,3)$. The $K$-label cut controls both retention of the
routing observation and the subsequent informative action choice.
\end{theorem}
\begin{proof}
For a fixed $i$, let
\[
                         p_i=\sum_s E(s\mid i)D(i\mid s).
\]
This is the probability of taking the informative action. On its
complement, all retained information is independent of $j$. Averaging
over the two values of $j$ makes the expected validation score zero
there, for any frozen event or randomized response. On the informative
part the score is at most $\TV(\mathrm{Ber}(1/2),\delta_j)=1/2$.
Thus the minimum over $j$ at a given $i$ is at most $p_i/2$.
A stop before sensing is also uninformative about $j$ and obeys the
same bound; it can be padded by an ignored sensing operation.

Set $L=\lceil m/K\rceil$. Since $p_i\le\max_sD(i\mid s)$,
a value $\min_i p_i>1/L$ would require every $i$ to have some row $s$
with $D(i\mid s)>1/L$. Each stochastic row has at most $L-1$ such
indices. But $K(L-1)<m$. This contradiction proves the upper bound.
Partition $[m]$ into at most $K$ fibers of size at most $L$, retain the
fiber label after $R$, and choose uniformly inside that fiber. On a
successful report freeze $\{1\}$ when $j=0$ and $\{0\}$ when $j=1$;
on $\perp$ freeze the empty event. The score for every $j$ is
$1/(2|\text{fiber}(i)|)$, proving equality. The event alphabet has three
labels. Storing the event and its candidate indicator requires at most
six labels, and the final difference has three. The chosen sensing
action was separately charged, proving the displayed profile.

For the convexified upper bound use the uniform prior on $(i,j)$.
For any one controller,
\[
 \frac1m\sum_i p_i
    =\frac1m\sum_{i,s}E(s\mid i)D(i\mid s)
    \le\frac1m\sum_{i,s}D(i\mid s)=\frac Km.
\]
The same upper bound holds for mixtures. For the matching mixture,
choose uniformly one of the $m$ cyclic intervals of $K$ indices.
For that interval hardcode a deterministic $K$-label encoder which
routes each included index to its own sensing action, and assigns other
indices arbitrarily. Each $i$ belongs to exactly $K$ intervals. A
successful action is followed by the complementary singleton event and
failure by the empty event, as above. Every fixed $(i,j)$ has mean score
$K/(2m)$. This proves \eqref{eq:routing-mixed}.
\end{proof}

For example, $V_{3,2}=1/4$ whereas $\overline V_{3,2}=1/3$.
The mixture just constructed can be implemented with a retained interval
selector and a local label, using at most $mK$ labels at the bottleneck.
It cannot be declared a $K$-label realization: when the displayed gap is
positive, \eqref{eq:routing-value} proves that no such realization exists.
The action buffer of size $m$ is not part of the $K$ coordinate and has
not been omitted from the profile. The theorem is a multicut resource
law, not an assertion that its uniform peak equals $K$.

The actions are statistically incomparable. $R$ cannot distinguish the
two values of $j$ at any fixed $i$, while action $a$ does distinguish
them at $i=a$. For $a\ne b$, action $a$ does not distinguish the pair
with $i=b$, while $b$ does. A parameter-independent garbling cannot
create a distinction absent from its input experiment.

\begin{corollary}[A strict gap with positive report probabilities]
\label{cor:routing-noise}
Replace each of the two training kernels by a mixture with weight
$\epsilon>0$ of the uniform law on its report alphabet. Let
$\Delta_{m,K}=K/(2m)-1/(2\lceil m/K\rceil)>0$.
For $\epsilon<\Delta_{m,K}/8$ the resulting full-support model still
has a strict gap between the genuine profile-constrained value and its
free-selector convexification.
\end{corollary}
\begin{proof}
Each training row changes by at most $\epsilon$ in total variation.
Couple corresponding rows until the first disagreement. For any fixed
controller or mixed controller the probability of such a disagreement
is at most $2\epsilon$, and its conditional reward lies in $[-1,1]$.
Each of the two optimal values changes by at most $4\epsilon$.
Their difference is therefore at least $\Delta_{m,K}-8\epsilon>0$.
All row arguments and register alphabets are unchanged.
\end{proof}

The distinction between mixed and behavioral strategies without perfect
recall is classical; it is not first asserted here. Contemporary
polynomial formulations over behavioral simplices include
\cite{ZSV}. Theorem~\ref{thm:routing-gap} evaluates the discrepancy in
this specific frozen-validation experiment for every $m,K$, with a
matching actual register and a noisy persistence statement. In
particular, a single prior followed by an unrestricted policy-tree
best response cannot be substituted for \eqref{eq:memory-dp}.
